% Handout-Vorlage zu den Lernvideos (gleiche Farben wie im Video).  Übersetzen:  pdflatex vorlage.tex
\documentclass[11pt,a4paper]{article}
\usepackage[utf8]{inputenc}
\usepackage[T1]{fontenc}
\usepackage[ngerman]{babel}
\usepackage{lmodern,amsmath,amssymb}
\usepackage[margin=2.2cm]{geometry}
\usepackage{xcolor,tikz,tcolorbox}
\usetikzlibrary{arrows.meta}

% Farben der Videos
\definecolor{vgold}{HTML}{E0A82E}
\definecolor{vcyan}{HTML}{1F9CC4}
\definecolor{vpink}{HTML}{D81B72}
\definecolor{vnavy}{HTML}{0D1430}

\newtcolorbox{merke}{colback=vcyan!6, colframe=vcyan, title=Merke, fonttitle=\bfseries}
\newtcolorbox{aufgabe}[1]{colback=vgold!8, colframe=vgold!90!black, title={Aufgabe #1}, fonttitle=\bfseries}

\title{\color{vnavy}\textbf{Titel des Lernvideos}\\[2pt]\large Handout}
\date{}
\begin{document}
\maketitle

\section*{Worum es geht}
Kurze Zusammenfassung des Videos in zwei, drei Sätzen.

\begin{merke}
Die eulersche Formel: \quad $e^{i\varphi} = \cos\varphi + i\sin\varphi$
\end{merke}

\begin{center}
\begin{tikzpicture}[scale=1.8, line width=0.8pt]
  \draw[-{Stealth}] (-1.3,0) -- (1.4,0) node[right] {$\mathrm{Re}$};
  \draw[-{Stealth}] (0,-1.3) -- (0,1.4) node[above] {$\mathrm{Im}$};
  \draw[vcyan] (0,0) circle (1);
  \draw[vgold, line width=1.4pt, -{Stealth}] (0,0) -- (40:1) node[above right] {$e^{i\varphi}$};
  \draw[dashed] (40:1) -- ({cos(40)},0) node[below] {$\cos\varphi$};
  \draw[dashed] (40:1) -- (0,{sin(40)}) node[left] {$\sin\varphi$};
  \draw (0.35,0) arc (0:40:0.35) node[midway, right] {$\varphi$};
\end{tikzpicture}
\end{center}

\begin{aufgabe}{1}
Zeige mit der Reihe $\displaystyle e^{x} = \sum_{n=0}^{\infty} \frac{x^n}{n!}$, dass $e^{i\pi} + 1 = 0$ gilt.
\end{aufgabe}

\end{document}
